\documentclass[11pt]{article}
\usepackage[a4paper,margin=25mm]{geometry}
\usepackage{amsmath,amssymb,booktabs,longtable,array}
\usepackage[colorlinks=true,urlcolor=blue,linkcolor=blue]{hyperref}
\title{Interleaving physical constraints with sewing:\\the $17+3$ and $16+4$ ten-loop MHV hexagon calculation}
\author{Computational technique report}
\date{Result completed 8 October; extended validation 9 October 2026 (Pacific time)}
\begin{document}
\maketitle
\begin{abstract}
The ten-loop weight-20 MHV hexagon symbol is obtained in the native
$17+3$ coordinates and converted exactly to the requested $16+4$
coordinates, within the explicitly stated conditional ansatz.
The native rational export has 246941 nonzero entries in 275246
coordinates; the equivalent shifted export has 491197 nonzero entries
in 518048 coordinates. An independent WXF reader confirms equality of
every coefficient after exact conversion. Separate denominator-bound
certificates verify all 2053259 original native joining equations and
all 4196956 original shifted joining equations over $\mathbb Q$.
All 44 reserved origin predictions pass exactly in both representations.
The native numerator and denominator heights are 45 and 315 bits; the
shifted heights are 44 and 231 bits. The shifted six-prime candidate was
validated against two subsequently computed primes. Native reconstruction
and its additional retrospective checks have separate provenance.
The method interleaves physical equations with integrability while retaining
symmetry through a common affine coordinate space. Measured changes to
allocation, OPE storage, MRK evaluation, scheduling and rational
reconstruction make the calculation feasible within a 50 GB decimal cap.
The double-final-entry bonus relations remain extrapolated assumptions;
finite sampled collinear, MRK and OPE checks do not constitute a proof of
all physical identities or unrestricted ten-loop uniqueness.
\end{abstract}

\section{The problem and the space in which it is solved}
At loop order $L$, the ordinary MHV symbol has weight $2L$. The requested
split at $L=10$ uses a weight-16 first-entry coproduct basis and a weight-4
last-entry coproduct basis. Their internal constraints are already imposed;
the remaining integrability and adjacency equations act at their joining
interface. Constructing the full interface matrix, or a large rational
kernel before imposing physical data, is the computational bottleneck.

The first-entry chain implements internal integrability, extended Steinmann
conditions and its specified one-term adjacency constraints. The generic MHV
last-entry chain implements internal integrability and full adjacency from
the fixed single-final-entry seed. The joining map uses the supplied full-CA
tensor. These are the input spaces actually used, rather than an unrestricted
weight-20 word ansatz.

The computation works directly in the parity-even $D_3$ invariant tensor
space. Its coordinates include all matching irreducible sectors:
\begin{center}
\begin{tabular}{lrrr}\toprule
Sector & First copies & Last copies & Invariant coordinates\\\midrule
$\mathrm{pe}\ A_1$ & 3445 & 28 & 96460\\
$\mathrm{pe}\ A_2$ & 2771 & 16 & 44336\\
$\mathrm{pe}\ E$   & 6213 & 44 & 273372\\
$\mathrm{po}\ A_1$ & 902 & 25 & 22550\\
$\mathrm{po}\ A_2$ & 873 & 14 & 12222\\
$\mathrm{po}\ E$   & 1772 & 39 & 69108\\\midrule
Total & & & 518048\\\bottomrule
\end{tabular}
\end{center}
The physical dimensions of the two bases are 23961 and 249. Their full
product has 5966289 coordinates. Passing to the invariant multiplicity
coordinates gives 518048 without discarding the $A_2\otimes A_2$,
$E\otimes E$, or odd-parity-pair contributions. The fixed Clebsch--Gordan
conventions and physical-copy-primitive-v1 normalization are retained.
The two components of $E$ are not independently rescaled.

\section{The alternative $17+3$ split and conversion to the requested coordinates}
The user also authorized the $17+3$ split. It moves one coproduct letter
from the last-entry basis into the first-entry basis, while keeping total
weight 20. The completed FEC17 basis has physical dimension 38472; LEC3 has
physical dimension 83. Their matching invariant sectors are
\begin{center}
\begin{tabular}{lrrr}\toprule
Sector & First copies & Last copies & Invariant coordinates\\\midrule
$\mathrm{pe}\ A_1$ & 5517 & 11 & 60687\\
$\mathrm{pe}\ A_2$ & 4449 & 4 & 17796\\
$\mathrm{pe}\ E$ & 9959 & 14 & 139426\\
$\mathrm{po}\ A_1$ & 1455 & 9 & 13095\\
$\mathrm{po}\ A_2$ & 1415 & 5 & 7075\\
$\mathrm{po}\ E$ & 2859 & 13 & 37167\\\midrule
Total & & & 275246\\\bottomrule
\end{tabular}
\end{center}
This is 46.8686\% fewer invariant coordinates than $16+4$. It does not
predict a corresponding runtime reduction: sparse elimination fill and
coefficient heights depend on the split. FEC17 generation took 1596.590
seconds, peaking at 44.411 GB; it has 171950436 compact nonzeros and
343-bit coefficient height. Its six blocks and completion manifest are
verified before use.

A completed $17+3$ finite-field symbol can be expressed in $16+4$
coordinates before rational reconstruction. The converter expands the first
recurrence, recouples the fixed CG channels, and solves for coordinates in
the target last-entry recurrence. Multiplication back through that target
recurrence must reproduce the expanded tensor exactly modulo the prime.
Each output cache binds both recurrence identities and the hash of its
source cache. A derived $16+4$ cache is not a native $16+4$ solve.

This route has passed a known-symbol control: the eight-loop $13+3$ solve
was unique modulo $p=1152921504606847009$, and its left shift to $12+4$
reproduced every coefficient of the established eight-loop result. The
$13+3$ solve took 110.391 seconds with a 2.227 GB peak; conversion took
0.406 seconds with a 0.160 GB peak. These figures establish the converter
and sewing conventions on the control, not a ten-loop completion time.
The completed ten-loop target has rational reconstruction and an independent
full $16+4$ joining certificate, detailed below. This check is necessary
because moving the split does not by itself establish equality of the two
constrained ansatz spaces.

\section{Why physical equations may be imposed before sewing is complete}
Let $V$ be this invariant coordinate space and $M$ the complete joining map.
The integrable symbols satisfy $Mc=0$. Suppose the current affine candidate
is written as
\begin{equation}\label{eq:affine}
 c=o+N\alpha,
\end{equation}
where the columns of $N$ span its surviving directions. A new block of
equations $Ac=b$ reduces to
\begin{equation}
 (AN)\alpha=b-Ao.
\end{equation}
If its solution is $\alpha=a+K\beta$, the next candidate is
\begin{equation}\label{eq:update}
 c=(o+Na)+(NK)\beta.
\end{equation}
Every earlier equation remains true for both the new origin and every new
direction. Homogeneous sewing has $b=0$; physical data can have $b\ne0$.
Consequently, retaining only homogeneous kernels after a physical fit would
lose essential information. The origin is stored as the final row of every
candidate, checkpoint, and native-prime cache.

This algebra permits any ordering of consistent linear blocks. If a
physical functional $P$ is initially defined on $\ker M$, its implemented
linear extension $\widetilde P$ to $V$ must agree with $P$ on $\ker M$.
Under that condition,
\begin{equation}
 \ker M\cap\{c:\widetilde Pc=b\}
 =\{c\in\ker M:Pc=b\}.
\end{equation}
An off-integrable candidate is therefore only a temporary linear-algebra
object. All remaining sewing equations must eventually be imposed. Early
physical restrictions are not a substitute for this final intersection.

Symmetry is compatible with this procedure because $o$, every column of
$N$, and all subsequent affine combinations stay in $V$. In the actual
storage convention, directions and origin are rows, so equation
\eqref{eq:update} is implemented as one common left multiplication on each
of the six irrep blocks. Giving different blocks independent solution axes
would break the common coefficient assignment; the implementation uses the
same affine transform for all six.

\section{Putting physical rows into the initial sewing reduction}
The scheduling requirement is to make all required physical equations, or
their lazy block generators, available alongside integrability and symmetry
from the start. Assembly and elimination may occur later. A costly block
$Ac=b$ evaluated on the current affine coordinates $c=o+Ny$ becomes
$ANy=b-Ao$; a deferred block must be regenerated or transformed consistently
after each coordinate change, retaining the common origin. Completion
accounting and final checks must prevent omission or stale projected rows.
This permits cost-aware block execution without requiring physics to be
solved first. The completed implementations below test specific schedules;
they do not establish a fully generic pending-block scheduler.

The preceding affine argument also permits physical constraints to enter
before discovery has constructed a small candidate. This is the user's
requested refinement for $17\otimes3$. V10 implements it using physical
rows in the original invariant coordinates, followed by a common augmented
reduction. The former schedule first reduced sewing to approximately 9000
directions and only then fitted physical data; that schedule can accumulate
substantial fill before physics has supplied any pivots.

The collinear and MRK projection already evolves separate first and last
basis images. Let $H_{bf\mu}(s,j)$ and $L_{bl\mu}(s,j)$ be their values for
sample $s$, irrep $b$, copy indices $f,l$, component $\mu$, and MRK
logarithmic degree $j$. A direct physical row has coefficient
\begin{equation}
 A_{s,q;(b,f,l)}=\sum_\mu g_\mu\sum_{j=0}^{q}
 H_{bf\mu}(s,j)L_{bl\mu}(s,q-j),
 \qquad g_0=1,\quad g_1=3.
\end{equation}
Collinear scalar evaluation has only degree zero. The component metric is
the existing fixed irrep convention. The MHV right-hand side retains the
existing factor $1/2$. The generator forms these coordinate images directly;
it does not project a matrix with 275246 independent identity solutions.
For the three-letter last side, collinear projection can instead keep its
boundary suffix words exact while sampling only the first side. Its
right-hand side uses those same suffix coordinates, including words outside
the last-entry image.

Introduce one shared scalar $t$ and reduce
\begin{equation}
 [M\ \ 0]\begin{pmatrix}c\\t\end{pmatrix}=0,
 \qquad [A_j\ \ -b_j]\begin{pmatrix}c\\t\end{pmatrix}=0.
\end{equation}
Thus sewing rows and physical rows belong to the same streaming reducer.
The extra coordinate is last and is not permuted. After the cheapest sewing
seed, one exact-tail collinear sample and one MRK sample are inserted. Fresh
batches are inserted between the next three sewing refinements. Physical
rows bypass the sewing row-density and sampling filters. They are sorted
by support within each batch, but their equations and right-hand sides are
retained exactly.

If $k_0$ is a kernel row with $t_0\ne0$, define $o=k_{0,c}/t_0$. A different
kernel row $(k_c,t_k)$ becomes the homogeneous direction $k_c-t_k o$.
Removing the $t$ coordinate yields one affine origin and a common set of
directions. If every kernel row has $t=0$, the affine system is inconsistent;
there is no acceptable physical origin. The temporary permutation is
inverted after this slicing. All six sectors use the same $t$, so symmetry
and a common coefficient assignment are preserved.

The direct-row probe compares every coefficient with the original identity
projection at a small control split, then contracts the rows into arbitrary
sparse directions and a nonzero origin. It checks both $Ao$ and $b-Ao$ for
scalar collinear, MRK and exact-tail collinear projection. The affine probe
checks a unique solution, a surviving direction, a nonzero-$t$ origin row
at a nonfinal position, permutation recovery and inconsistent RHS rejection.
Both probes passed. The full eight-loop V10 control also passed every
remaining sewing row and fresh physical check; its unique coefficient
payload agrees with V9 exactly modulo the test prime. Early physical
batches added 94 pivots. Discovery left 1869 directions rather than 1963;
after the subsequent collinear fit, 66 remained rather than 112. Both runs
then reached 14 after MRK and zero after the covered OPE fit. V10 took
124.808 seconds with a sampled 2.316 GB peak, compared with 112.339 seconds
and 2.197 GB for V9. Its candidate had 19420234 nonzeros versus 19301649.
Thus earlier physical information reduced the remaining dimension in this
control, but did not demonstrate a runtime or fill improvement.
The first ten-loop V10 trial was frozen during validation and then retired
after the V13 nine-loop and V15 eight-loop controls passed. Its input files,
and logs are preserved; obsolete temporary working storage was subsequently
removed under the authorized scratch cleanup. Its last reported rank was 248552, with
3331074850 retained nonzeros. The live field entries alone require about
39.97 GB before input arrays and allocation overhead. Repeated paging made
low processor utilization a symptom of memory traffic, rather than evidence
that more threads would resolve the bottleneck. Neither that trial nor its
V15 replacement completed a native ten-loop prime. Both used a 48 GB
pressure threshold and a 50 GB hard limit.

This ordering offers earlier pivots, not a theorem that every new row will
be sparse. Generic physical rows can themselves create fill. Bounded batches
make that cost measurable without materializing a full physical kernel.
All remaining adaptive collinear, MRK and covered OPE equations are still
imposed on the discovered candidate; complete residual sewing and fresh
physical checks remain mandatory.

\section{Direct leading OPE constraints and measured scheduling}
The OPE conditions can enter the same augmented reduction. To avoid evolving
275246 separate unit candidates, the direct implementation uses the
transpose of the existing finite jet integration map. If $J_{\ell,w}$
integrates a jet against letter $\ell$ at weight $w$, then
\begin{equation}
 \langle e,J_{\ell,w}h\rangle
 =\langle J_{\ell,w}^{T}e,h\rangle.
\end{equation}
The first-entry jets are evolved once for each helicity evaluation. A chosen
OPE coefficient supplies a sparse endpoint covector; the last-entry chain
propagates that covector backwards through weights $20,19,18$. Contracting
the resulting last-side cotangent against the first-side jet gives one row
in the original invariant coordinates, with the same component metric and
single affine normalization as above.

The transpose reverses all three stages of the existing integration:
origin logarithmic integration, the triangular logarithmic recursion, and
kernel convolution. It is tested by arbitrary primal--dual dot products for
all letters and weights in the small control. The resulting coordinate rows
are also multiplied into arbitrary sparse directions and a nonzero origin,
and compared with the unchanged grade-12 projection. Both checks pass.

Only existing solve-role $T S$ and $T^2 S^2$ coefficients are selected for
early use. Grade four suffices for these targets because every ordinary
kernel term has nonnegative total power increment, while logarithmic
kernel terms do not lower the retained grading; the implementation checks
this filtration condition. The same three evaluations $F=2,3,4$ eliminate
the unknown $T^2F^0$ term. No held-out coefficient is used for fitting and
no unavailable coefficient is set to zero. At ten loops one grade-four
first-side jet occupies approximately 2.04 GB, so it is released between
helicity evaluations instead of retaining all three beside the reducer.

The V11 experiment continues disjoint collinear and MRK sample streams
after the four light rounds, activates larger batches once augmented
nullity reaches $\min(50000,n/5)$, and introduces leading OPE after two
initial cycles. Each subsequent callback has a soft time budget, checked
between batches. Rank gained, elapsed time and retained storage are logged
to compare scheduling choices. Stagnation thresholds and sample caps guide
discovery only; the later full sewing sweep and fresh physical checks
remain the acceptance conditions. More early equations need not reduce
runtime, because their evaluation and the fill they create also cost work.

The associated allocator change makes shrinking a sparse buffer return its
unused suffix to the mapped arena. Adjacent free extents coalesce, and only
whole pages inside free storage may be hole-punched; partial pages shared
with live arrays are retained. The allocator probe checks live neighbours,
tail reuse, suffix trimming and forced paging. The reducer comparison also
preserves exact coefficients and pivot order against V8 at two primes.
These engineering checks establish storage correctness; only completed
amplitude controls can establish the benefit of a complete schedule.

The completed V11 $13+3$ eight-loop control reproduces the V10 coefficient
payload exactly modulo $1152921504606847009$, with zero surviving directions
and all final sewing and physical checks passing. Early leading OPE adds
255 independent pivots in this control. Its discovery candidate has
14618598 nonzeros, compared with V10's 19420234. Nevertheless the complete
run takes 265.720 seconds and peaks at 2.587 GB, compared with 124.808 seconds
and 2.316 GB for V10. The aggressive schedule therefore improves candidate
fill but fails the end-to-end time and memory comparison on this control;
it is not selected as the final ten-loop strategy on that evidence.

\section{A conditional restriction before dense elimination}
The user permits reasonable additional assumptions, provided that any unique
candidate is subsequently tested against independent physical information.
The next experiment therefore uses the empirical MHV double-final-entry
relations of Dixon and Liu, equations (76)--(79) in \cite{dixonliu}. These
are additional to the ordinary double-final-entry $\bar Q$ relations:
the latter already vanish on the generic 25-dimensional pair space and
would supply no new constraints. The bonus relations are treated explicitly
as an extrapolation to ten loops, not as an all-orders theorem.

The relations are published in the alphabet
$(u,v,w,1-u,1-v,1-w,y_u,y_v,y_w)$. Let $B$ express the working alphabet's
logarithmic differentials in this alphabet. A functional on the last two
coefficients transforms by $R\mapsto BRB^T$. The implementation calibrates
this dual transformation against both the ordinary single-entry and the
ordinary double-entry conditions, before admitting the bonus equations.
It expands only the small last-entry chains through weight four and applies
the full $D_3$ orbit of each relation to the final two slots. The exact
copy kernels give
\begin{center}
\begin{tabular}{lrr}\toprule
Last-entry weight & Generic physical dimension & With double bonus\\\midrule
2 & 25 & 21\\
3 & 83 & 65\\
4 & 249 & 184\\\bottomrule
\end{tabular}
\end{center}
For the three-letter side, the reduced copy multiplicities, in the six
sector order used above, are $(9,3,11,8,3,10)$. Hence the conditional
weight-20 invariant ansatz has 217024 coordinates at $17+3$, compared with
275246 before these relations. At $16+4$ the corresponding count is 380474,
compared with 518048. These counts precede the joining constraints and
do not assert the dimension of the physical solution space.

The audit verifies the exact rational copy kernels, their right inverses,
and coordinate round trips. It also contracts every bonus constraint with
the completed eight-loop symbol at the first native prime, and with the
completed nine-loop symbol at two native primes. Every residual is zero.
The common copy map acts on both components of each $E$ sector, preserving
the existing symmetry and normalization conventions.

In the solver experiment, independent bonus equations are inserted as short
homogeneous rows in the original invariant coordinates after the cheap
sewing seed and before dense refinement. This avoids changing either the
generic last-entry files or the subsequent $17+3\to16+4$ conversion. The
candidate descriptor records the empirical assumption. A unique result in
this restricted ansatz would establish uniqueness only within that ansatz;
original sewing and held-out physical tests provide independent evidence
for the candidate, not a proof of unrestricted ten-loop uniqueness.

The conditional V12 eight-loop control passes every original sewing and
fresh physical check and reproduces the known coefficient payload exactly.
Its 8228 short bonus rows add 6145 pivots after the seed in about 0.04 seconds.
Discovery has 7585975 nonzeros, and the complete run takes 160.058 seconds
with a 2.367 GB sampled peak. This improves on the aggressive V11 schedule,
but still costs more time than V10.

V13 therefore limits each physical callback to one cycle and adds callbacks
at completed elimination-batch boundaries when 120 seconds have elapsed
inside a long sewing sweep. The row-generation workers and reduction workers
have joined before such a callback. A previously generated row is reduced
against the updated basis when consumed; the coordinate permutation remains
fixed. A targeted affine probe exercises this mid-sweep path with a nonzero
right-hand side and a nontrivial variable permutation. It passes, and the
complete eight-loop control again reproduces the known symbol, with all
final checks passing. Discovery has 7923078 nonzeros; runtime is 129.512
seconds and sampled peak memory is 2.319 GB. This is close to V10's complete
runtime while substantially reducing discovery fill.

The larger V13 nine-loop $15+3$ control has now also passed. Discovery takes
297.101 seconds and returns 5962 directions with 48535187 nonzeros; collinear
equations reduce this to 564 directions. After the remaining physical fit,
complete sewing reduces the candidate to a single point, and all fresh
physical checks pass. The complete native control takes 1474.609 seconds
with a 7.755 GB sampled peak. Conversion to $14+4$ takes 5.222 seconds and
reproduces every coefficient of the independently completed nine-loop
symbol modulo $1152921504606847009$. Its wall time includes overlapping
compilations, a FEC17 memory probe and the later eight-loop controls; it
must not be treated as an isolated throughput benchmark.

\section{Recovering memory retained by integer caches}
A real FEC17 input probe identifies an additional engineering bottleneck.
After the temporary rational recurrence has been destroyed, its finite-field
cache contains 171950436 entries and occupies about 2.06 GB, but the process
still retains about 10.84 GB. Ordinary allocator collection returns only a
few megabytes: FLINT still owns reusable multiprecision integer objects in
thread-local caches, so their limbs are not yet free allocator pages.

V14 calls FLINT's cache cleanup at joined phase boundaries, after releasing
the temporary rational input. One task per worker waits at a barrier before
cleanup; this guarantees coverage of every worker, unlike submitting several
short tasks that may all run on one worker. The main thread is cleaned as
well. The global master cleanup routine is not used during the solve.
The implementation follows the installed FLINT 3.2.2 cache ownership rules
and its documented thread-local cleanup mechanism \cite{flintmemory}.

In the FEC17 probe this operation reduces resident memory from
10.8467 GB to 2.5913 GB in 1.37 seconds. All six finite-field block hashes
remain identical to the previous probe, every live rational sentinel is
preserved, and subsequent large-rational arithmetic succeeds. Cleanup
releases unused cached objects; it does not change a mathematical basis or
discard a retained recurrence. The complete V14 eight-loop solver reproduces
every known coefficient and passes its original sewing and physical checks.
It takes 156.210 seconds with a 2.322 GB sampled peak while the independent
nine-loop control is also running; this is a correctness control, not an
isolated runtime comparison.

The finite-field input cache is also sized to retain the whole needed chain:
the FEC17 residue alone needs 2064226160 bytes, so a nominal 2 GiB cache
cannot retain it together with earlier weights. Repeatedly decoding the
same rational input would negate the benefit of bounded physical batches.
The proposed ten-loop configuration uses a 4 GiB recurrence cache. The
square logarithmic layout would require approximately 5.29 GB for each of
three grade-12 FEC17 jets. The next optimization reduces this requirement
before assigning the late OPE cache. Early OPE continues to release its
grade-four jets between evaluations while the elimination basis is live.

\section{Triangular logarithmic storage for OPE jets}
The OPE evaluation integrates the symbol recurrence into truncated series in
$T$ and $S$, with coefficients polynomial in $\log T$ and $\log S$.
Let $W$ denote the final weight, and write a logarithmic monomial as
$(\log T)^a(\log S)^b$. At intermediate weight $w$, the integration recurrence
produces only $a+b\leq w$. Indeed, the weight-zero scalar has no logarithms,
and each integration either preserves the logarithmic degree or raises it
by at most one. Thus the final representation needs only $a+b\leq W$.
Storing all $0\leq a,b\leq W$ leaves a triangle of entries identically zero.

For each retained power pair of $T$ and $S$, V15 packs the logarithmic index
as
\[
 i(a,b)=a(W+1)-\frac{a(a-1)}{2}+b,
 \qquad 0\leq a\leq W,\quad 0\leq b\leq W-a.
\]
The number of slots is $(W+1)(W+2)/2$ instead of $(W+1)^2$.
At $W=20$ this changes 441 slots to 231, a reduction of $47.62\%$.
The forward and transposed integration maps both omit neighbors whose total
logarithmic degree exceeds the current weight. They retain the same field
coefficients, truncation order, target rows and normalization.

This change is checked against a separately retained square implementation,
not against another use of the new index. The probe visits all nine letters,
every weight from 1 through 20, truncation grades 4 and 12, and two primes.
It compares every retained output, verifies that discarded outputs are zero,
and checks the forward--adjoint pairing. All 2160 cases pass, including
inputs in logarithmic slots that are unreachable from the weight-zero seed.
The complete V15 eight-loop solve then passes its original joining and fresh
physical checks and reproduces the known coefficient payload exactly.
Its sampled peak is 1.492 GB, compared with 2.319 GB for V13. Its wall time
is 134.307 seconds while the nine-loop control is also active; no isolated
speedup is inferred from those overlapping measurements.

For the ten-loop grade-12 calculation, each FEC17 head now has stride 9009
and occupies approximately 2.773 GB. Three such heads fit in an 8 GiB cache,
so later physical checks can reuse all three evaluations without repeating
the long recurrence integration. At grade 4, the early head stride changes
from 6615 to 3465, approximately halving its temporary workspace as well.
This is an exact representation change and introduces no further physical
assumption beyond the conditional final-entry restriction stated above.

\section{Reusing common factors in the MRK right-hand side}
Ten-loop profiling identifies a bottleneck that memory reclamation does not
address. Each of the first single-sample MRK right-hand sides takes roughly
140 seconds, while its head and tail recurrence projections take only a few
seconds. A 30-second resource sample records approximately one busy core
and no memory pressure. The reference generator repeatedly evaluates the
same shuffle products for different terms in the physical formula.

The replacement shares those common subexpressions. For slot $i$, write its
three sampled field values as $(a_i,b_i,c_i)$ and let $x$ track logarithmic
degree. Define
\[
 \tau_i(x)=\tfrac12 b_i-c_i+a_i x,
 \qquad Q_{ij}(x)=-2\tau_i(x)\tau_j(x)-\tfrac12 b_i b_j.
\]
These are precisely the existing one-letter logarithmic and symmetric
one-loop two-letter factors. Let $H_{p,d}(I)$ denote the evaluation on the
slot subset $I$ of $p$ shuffled copies of $Q$, divided by $p!$, shuffled
with $d$ copies of $\tau$. It vanishes unless $|I|=2p+d$, with
$H_{0,0}(\varnothing)=1$. Choosing the least slot $i\in I$ gives
\[
 H_{p,d}(I)=d\tau_i H_{p,d-1}(I\setminus\{i\})
 +\sum_{j\in I\setminus\{i\}} Q_{ij}
 H_{p-1,d}(I\setminus\{i,j\}),
\]
where terms with negative indices are zero. The first term counts the $d$
identical logarithmic factors that could occupy $i$. The second pairs $i$
with its partner; dividing by $p!$ removes the ordering of the identical
one-loop factors. Symmetry of $Q_{ij}$ makes this a sum over unordered
pairings with the correct shuffle multiplicity.

For a fixed slot set, $d=|I|-2p$, so the implementation memoizes only
$(p,I)$. It then contracts each input MRK symbol on its complementary ordered
subset. Binary words and their coefficients modulo the current prime are
prepared once per evaluation. No symbol coefficient or sample is changed.
The final derivative and half-phase contributions combine into $\tau_i$,
using the same expression above. All intermediate polynomials are truncated
at the requested nonnegative degree, which commutes with these polynomial
sums and products. The external MHV factor of one half remains unchanged.

There are 32 exact comparisons against the retained generator, covering
loop orders 1--10, small and production-size primes, different sample
indices, and truncated as well as full logarithmic degree. Every returned
coefficient agrees. For the two ten-loop cases the isolated reference/new
times are 138.420/21.0852 seconds and 143.734/21.0271 seconds, respectively:
a factor of approximately 6.6--6.8 for this kernel. The larger comparison
process peaks at 0.420 GB. These measurements are made sequentially with
the numerical pilot frozen; they do not establish a complete-solver speedup.

V16 integrates this generator and precomputes four RHS samples when a
single-sample callback is requested. Only the requested equations are
admitted, in their original order; the other RHS values are retained for
later callbacks. Cache keys include prime, seed, loop order, logarithmic
degree and sample index. At most 128 samples are retained, and larger
batches bypass retention. Generation uses the existing worker pool from a
joined outer call, avoiding nested waits. Cache and complete known-symbol
controls have now passed. The cache probe checks 1848 cases across prime,
seed, degree and sample changes, including cache hits and large uncached
batches. With the pilot frozen and compilation finished, sequential complete
eight-loop runs take 103.267 seconds for V15 with 8 threads, 76.712 seconds
for V16 with 8 threads, and 60.730 seconds for V16 with 16 threads. Every
coefficient agrees with the independently known result. Thus V16 reduces
this complete-control runtime by 25.7 percent at equal thread count, and
16 threads reduce its runtime a further 20.8 percent. These are measured
control results, not a claimed ten-loop speedup or global optimum.

V16 with 16 threads was consequently selected for the ten-loop first-prime
pilot. The unfinished V15 trial was retired with its inputs and evidence
preserved. V16 has now completed one native prime and its conversion;
the scheduling review below selected V17 for the subsequent primes.

\subsection{Preserving unfinished physical work across deadlines}
The full-scale pilot exposes a scheduling defect that the faster controls
did not trigger. A four-sample collinear batch costs about 19--21 seconds
at the relevant ten-loop discovery dimension. Six or seven such batches
consume the 120-second callback budget before the intended eight-batch
collinear segment is complete. V16 restarts that segment at the next
callback and increments its cycle count even when MRK was skipped. Repeated
callbacks can therefore postpone both MRK and the due low-grade OPE block
indefinitely. In this pilot, discovery admitted 71 collinear batches but
only the initial four light MRK samples; no low-grade OPE block was admitted.
The subsequent full physical fits still ran, so this finding concerns
the timing of constraints rather than removal of the final checks.

The separate V17 correction retains a cursor counting collinear batches
already completed in the current cycle. Once their quota is reached, the
next eligible work is MRK; a new cycle starts only after that MRK batch is
admitted or its stream is exhausted. OPE becomes the first eligible block
once two cycles are complete, or both sampling streams have finished.
This also preserves an OPE obligation when the preceding MRK batch exceeds
the remaining callback budget. Each admitted collinear batch advances a
bounded cursor, so repeated eligible callbacks cannot repeatedly start
the same segment. The argument assumes callbacks continue and the target
dimension has not already been reached. A single indivisible block may
still exceed the nominal time budget.

The correction changes admission order and persistence of work, while
retaining the original affine equations, right-hand sides, symmetry
conventions and final checks. Its 5184 scheduler controls vary deadlines,
operator costs, stream exhaustion and early target attainment, checking
that the operator sequence is preserved without replay. The full known
eight-loop coefficient payload also agrees exactly and its final checks
pass. That control ran alongside the productive pilot, so its timing is
not an isolated speed comparison. The serialized nine-loop regression has
also passed: every native $15+3$ coefficient and every converted $14+4$
coefficient agrees with the independent completed controls. Its native
calculation took 576.942 seconds with sampled peak RSS 6.442 GB; conversion
took 4.821 seconds. These controls support the scheduling change, while its
effect on total ten-loop runtime is assessed in the completed measurements below.

To use V17 without discarding the completed ten-loop prime, a separate
output chain imports the V16 native and converted caches with an explicit
provenance receipt. The physical input hashes, recurrence basis, affine
origin convention, coefficient order and every residue remain identical.
The native descriptor changes only the algorithm, admission schedule and
executable hash. The cache header's eight-byte identity fingerprint is
rebound to that descriptor; the shifted identity is rebound to its new
parent identity. The original files and hashes remain preserved, and the
receipt identifies V16 as the producer of the imported prime. This is a
validated import, not another native solve or another sewing conversion.

The import was checked against the independently computed V17 eight-loop
native cache byte for byte. A V17 cache-resume check recomputed the input
identity and reused that cache without solving the prime. Its shifted
import also agrees byte for byte with an independent conversion of the
imported native cache. Six negative controls reject checksum corruption,
an incorrect identity, multiple affine rows, incorrect dimensions,
truncation and noncanonical residues. The supervisor pins the import
receipt, original evidence and validation hashes, and rechecks the full
source-to-import coefficient equality before final certification. The
final certificate records the original producer of each prime alongside
the independent rational sewing and origin checks.

\section{Factored sewing and staged discovery}
The first and last recurrence maps and the fixed joining tensor supply the
interface equations through contractions. The full product matrix is not
materialized. Left contraction gives 1232602 rows at this split; right
contraction gives 1925885. Both have 518048 invariant columns and enforce
the same complete joining condition after the required residual checks.

Discovery admits cheap sparse rows first, then progressively denser rows.
Exact singleton and doubleton relations support substitutions. Nested,
deterministic row samples guide the schedule. Sampling cannot certify a
symbol: the final acceptance step visits every unsampled joining equation.
The scheduling target is 12000 directions, not a prediction of the physical
dimension.

The right-contracted experiment temporarily permutes variables within each
irrep block. If that block has $m$ first copies and $n$ last copies, the
original coordinate $fn+l$ is sent to $lm+f$. Writing $c'=\Pi c$, discovery
uses $M\Pi^{-1}$. The candidate is mapped back by $\Pi^{-1}$ before physical
projection, checkpointing, residual sewing and reconstruction. The map is
checked to be bijective. It changes sparse elimination fill, not the ansatz,
the symmetry condition or coefficient normalization.

Once discovery supplies a manageable candidate, collinear, multi-Regge
(MRK), and covered operator-product-expansion (OPE) equations are imposed
in affine coordinates. The remaining joining map is then contracted into
that smaller candidate. A rank-changing residual block triggers an affine
restriction and a replay from the beginning. A complete sweep with no
remaining residual is required. Fresh physical checks follow on the fully
integrable object.

\section{Physical input coverage and the limits of the checks}
The ten-loop main and extended MHV OPE libraries have complete, hash-bound
generation receipts. The current run uses the main library, with grade 12,
solve power 8, and check power 10. Its T1 and T2 coverage supplies 1824 solve
equations and 456 held-out coefficients. The extraction combines evaluations
at $F=2,3,4$ to remove the unknown T2 $F^0$ contribution. An unsupported
coefficient is never replaced by zero. MRK inputs are available through ten
loops.

Physical constraints are applied using the supplied ordered projection
prescription. Collinear and MRK validation uses points separate from the
training stream. OPE validation uses its specified held-out powers. These
are concrete checks of the implemented conditions; finite sampled checks
alone do not prove an identity at every kinematic point. Full sewing has a
separate exact rational certificate, described below. Nor does an ordinary
symbol determine function-level constants invisible to the symbol map.

\section{Storage changes and what the controls establish}
Large intermediate objects must be shortened in lifetime as well as kept
sparse. Temporary rational recurrence matrices are released after the
finite-field joining operator has taken ownership of its required copies.
Kernel extraction consumes reduced pivot rows as it transfers their
coefficients to the kernel. Storage for trivial field elements is allocated
without eagerly clearing a second full array, and all live entries are
assigned before use. The free-variable sign convention remains the
original one: a private free coordinate is $-1$, with pivot coefficients
given by the corresponding reduced-row coefficients.

Physical projection transfers a packed candidate into its irrep blocks row
by row, releasing the input rows immediately. After fitting, the common
affine transform multiplies one block at a time; the old block is released
after its product is packed. This avoids retaining both a full packed
candidate and a complete second block copy.

Private intermediate checkpoints use sparse rows and include the affine
origin. For $r$ rows and $z$ nonzeros, their size is $80+8r+16z$ bytes,
compared with $72+8rn$ bytes for the conventional dense payload at $n$
coordinates. Writes are streamed back and compared exactly with the live
matrix; stage identity and SHA256 protect later reuse. Final native-prime
caches retain the package's original format. Wrong split identities,
damaged bytes, chunk boundaries, and nonzero origins are covered by the
checkpoint probe.

These changes and the variable permutation still left excessive discovery
fill. The V5 right-contracted trial stopped at 34.9 GB RSS when the system
reserve was exhausted. V6 reached a 43.1 GB peak; its last reported rank was
296591 with 2125869115 retained nonzeros. Neither produced a completed
native ten-loop prime. The retained logs distinguish last reported rank
from the rank at the actual memory-stop instant.

V7 therefore maps SparseRREF arrays of at least 4096 bytes into a fresh
private scratch file. Small arrays and GMP/FLINT limbs keep their ordinary
allocator. Shared file pages preserve exact bytes when residency is
released; later accesses reload those bytes. The allocator synchronizes
pages before requesting reduced residency. This is temporary working
storage, not an imported solution or a completed-prime cache.

The production mapping reserves 192 GiB of sparse-file virtual address
space under a 256 GiB virtual limit. Actual process-tree RSS is guarded at
exactly 50000000000 bytes (46.5661287308 GiB), with a 3 GiB system-memory
reserve. V15 raises the selected residency threshold from the previous
32 GiB to 40 GiB after reducing retained input and OPE storage. Input caches and ordinary
allocations also consume RAM, so the threshold is not a hard total-memory
bound. Sixteen internal threads operate on one native prime at a time.

The allocator probe passes growth, shrinking, heap-to-file transition,
concurrent-buffer and forced-eviction checks, including rational arrays
whose numerators and denominators exceed the small-integer representation.
More decisively, V7 reproduced the known eight-loop symbol exactly modulo
$p=1152921504606847009$, with all final sewing and physical checks passing.
That control exercised 28 evictions with a deliberately low 1 GiB
residency threshold. It took 148.675 seconds and peaked at 2.203 GB RSS.
Its final invariant coefficient payload has 65344 nonzeros. This validates
the storage and ordering changes on a known calculation; it is not a
ten-loop result.

\section{Reconstruction and an independent exact sewing certificate}
Only the final affine product is reconstructed over $\mathbb Q$.
Native good primes are accumulated by CRT and rational reconstruction;
two independent validation primes check the reconstructed coefficients.
Large intermediate presentation kernels are not reconstructed.

An independent checker contracts the reconstructed coefficients directly
through the original first recurrence, last recurrence and CG joining
tensor, rather than through the factored discovery reduction. Let
$\widehat F,\widehat L,\widehat T,\widehat c$ denote these four objects with
their respective common denominators cleared. Define
\begin{equation}
 B=\|\widehat F\|_1\,\|\widehat L\|_1\,
   \|\widehat T\|_1\,\|\widehat c\|_\infty.
\end{equation}
Every resulting integer residual has absolute value at most $B$. Good
certificate primes divide none of the cleared denominators. If every
residual vanishes modulo distinct primes whose product $Q$ exceeds $B$,
each residual is an integer multiple of $Q$ with magnitude less than $Q$,
and hence is exactly zero. The checker also compares the rational result
with all eight shifted caches and verifies each link to its native cache. Additional primes used solely for
residual certification are explicitly distinguished from native symbol
solves and never written as solve caches.

\subsection{Asymmetric reconstruction for denominator-heavy coordinates}
The production reconstruction used equal numerator and denominator bounds.
That is inefficient for this basis: the final largest numerator has only
44 bits, while the largest denominator has 231 bits. Six roughly 60-bit
fitting primes give a 361-bit CRT modulus $M$. A symmetric bound of roughly
$\sqrt{M/2}$ cannot accommodate that denominator, even though the same
residues contain enough information for unequal bounds.

The separate reconstruction experiment fixes
\[
 N=2^{48}-1,\qquad D=\left\lfloor\frac{M-1}{2N}\right\rfloor,
 \qquad 2ND<M.
\]
The bound $D$ has 312 bits; it is calculated from $M$, not set independently
to $2^{312}$. For a residue $r$, the Euclidean algorithm on $(M,r)$ tracks
$r_i\equiv t_i r\pmod M$. It stops when $|r_i|\leq N$, adjusts the sign so
that $t_i>0$, and requires $t_i\leq D$, coprimality and the original
congruence. If two reduced fractions within these bounds represented the
same residue, their cross-difference would be a multiple of $M$ of absolute
value at most $2ND<M$; it must therefore vanish. This proves uniqueness
inside the selected coefficient bounds, not that an unknown physical
answer necessarily satisfies them.

The implementation passed 2682 exhaustive small-modulus cases. Diagnostic
bound comparisons used existing residues and were treated as fitting
information. The final six-prime vector and bounds were then frozen before
native prime 7 was committed, with primes 7 and 8 reserved for future
validation. Reconstruction took 7.200 seconds. Every one of the 518048
entries, including zeros, agrees with both later primes; neither check
prime divides any candidate denominator. No coefficient or bound was tuned
after these checks. The two checks took 0.540 seconds. The six fitting
residues themselves were never counted as independent validation.

The rational vector was exported into a new directory rather than inserted
into the active producer's chain. A separate format-only exporter reads
canonical rationals in the fixed six-block order, writes reduced WXF
matrices, and reads them back exactly. An independent Wolfram WXF decoder
provided the known eight- and nine-loop controls: all 69988 and 195155
coordinates pass, and every exported block and metadata file is byte-identical
to its respective reference. Six malformed-input cases are rejected before
creating an output. Wolfram then independently checked all 518048 entries
of the ten-loop export. No basis change, normalization change, presentation
RREF or coefficient filtering is performed.

The new manifest explicitly records the external reconstruction and exporter.
The shift-map identity, basis identities, eight committed residue payloads,
and their source hashes are copied unchanged, retaining the original native
parent. Prime 1 remains attributed to V16 through its validated import;
primes 2--8 were computed by V17. The completion seal means the rational
object is written and hash-consistent; its physical acceptance follows only
from the additional checks below.

For this object the independent original joining operator has 4196956 rows.
Its denominator-cleared residual bound has 514 bits. All eight available
primes give zero residual; one further cheap residual evaluation increases
the certificate modulus to 541 bits, strictly exceeding the bound.
Consequently every original joining residual vanishes exactly over
$\mathbb Q$. This ninth certificate prime is not a ninth native symbol solve.
The unchanged certificate completed in 67.667 seconds with a sampled
4.398 GB peak RSS. It also verified all eight QQ-to-shifted-cache congruences
and all source-hash links.

\section{An independent origin prediction}
The origin check was held out of the fit and has now passed exactly. At ordinary
symbol level the all-order origin prediction of Dixon and Liu, equation
(42), gives the exponential generated by the one-loop origin term
\cite{dixonliu}. Since that term is
$-\tfrac12\sum_i\log^2u_i$ at the relevant logarithmic order, the coefficient
of a word with counts $n_i$ at loop order $L$ is zero if any count is odd,
and otherwise equals
\begin{equation}
 (-1)^L\prod_i\frac{n_i!}{2^{n_i/2}(n_i/2)!}
 =(-1)^L\prod_i(n_i-1)!! .
\end{equation}
Here $(-1)!!=1$. The factorials follow from the ordinary-symbol convention
for powers of logarithms; omitting them would give a different
normalization. At ten loops, counts $(20,0,0)$ predict 654729075, while
$(8,6,6)$ predict 23625. Pure-log integrability and $D_3$ symmetry reduce the
checks to 44 count partitions at weight 20. This prediction concerns the
origin limit and is not an identification with coefficients on the
$(u,1,1)$ line.

\section{Completed native calculations and reproducibility}
On 7 October the V15 first-prime trial identified the serial MRK generator
as a substantial bottleneck. After the exact and complete-solver V16
controls described above passed, the unfinished V15 trial was retired and
its inputs and evidence preserved. The V16 first-prime pilot has completed. It combines
the bounded physical schedule and conditional final-entry rows with integer
cache release, triangular OPE storage, shared MRK shuffle evaluation and
bounded RHS prefetch.

The pilot retained one native prime and its $16+4$ conversion, and stopped
before further primes for a review of full-scale timings and elimination
fill. The native prime is $1152921504606847009$. Its finite-field dimension
chain is
\begin{center}
\begin{tabular}{lr}
\toprule
Stage & Remaining affine directions\\
\midrule
Discovery candidate & 22489\\
Collinear fit & 3056\\
MRK fit & 3041\\
Covered OPE fit & 2877\\
Complete original sewing & 0\\
\bottomrule
\end{tabular}
\end{center}
The factored $17+3$ joining operator has 989800 rows; the finishing routine
made 35 passes including the final complete zero-residual sweep. Fresh
collinear, MRK and OPE checks then passed on the unique representative.
The native result contains 246941 nonzero coefficients in 275246 invariant
coordinates. It took 12711.891 seconds, about 3 hours 32 minutes, with
sampled peak RSS 34.628 GB and no resource-guard stop. The subsequent
$16+4$ conversion took 64.061 seconds including the reconstruction attempt,
peaked at 12.017 GB, and passed the exact finite-field shift checks. It has
491197 nonzeros in 518048 coordinates. The native and converted cache
hashes and their source linkage are committed in the V16 runner ledger.

After the first prime, the pilot supervisor exited in its expected
cache-only state. After the nine-loop scheduling regression and
explicit-import controls passed, V17 started at prime index 2 with a budget
of 12 primes, retaining the completed first prime and its conversion. It
holds both supervisor locks to prevent a concurrent V16 or V17 solve.
The hourly monitor followed this chain through reconstruction and exact
certification, including recovery after the host reboot. Correctness controls alone do not
establish practical optimization. The thread choice is supported by the
isolated eight-loop comparison; remaining ten-loop kernels and alternative
elimination orderings have not all been compared. Optimization continues
from measured bottlenecks while validated inputs and results are preserved.
The reusable decisions and unsuccessful alternatives are recorded in the
repository optimization lessons, referenced by its agent instructions.

The supervisor is configured to run one native prime at a time with 16 internal threads,
a 4 GiB recurrence cache and an 8 GiB late OPE head cache. After each
successful prime it converts the candidate to $16+4$ and attempts rational
reconstruction, reserving two independent validation primes. Its configured final stages are the independent exact original
$16+4$ sewing certificate and all 44 reserved origin tests. The faster
asymmetric reconstruction used a separately validated exporter and
controller for these same unchanged checkers, as detailed below. Completion is
reported only within the conditional ansatz and the stated physical
validation scope. The first-prime runtime established feasibility within the memory cap.
The completed reconstruction used six fitting primes and two fresh
validation primes; the completed native solves totaled 28.793 hours,
excluding stopped trials, input generation, downtime and certification.

The first full V17 calculation has now completed at the second prime,
$p=1152921504606847067$. Its native runtime is 12627.650 seconds, about
3 hours 30 minutes, with sampled peak RSS 38.581 GB. The verified $16+4$
conversion took 64.129 seconds and reused the already converted first
prime. All 989800 factored joining rows pass, with 35 finishing passes
and zero remaining directions; fresh collinear, MRK and covered OPE checks
also pass. The final native and shifted nonzero counts are respectively
246941 and 491197. Every completed cache hash and source link, and all five
sparse checkpoint identities and hashes, were rechecked independently of
the running supervisor. The first prime retains its explicit V16 producer
provenance. At that stage the supervisor continued with prime index 3: two completed
primes were insufficient for a protocol reserving two additional primes
to validate a CRT candidate. The final eight-prime outcome is given below.

The completed comparison explains why early physical rank gain alone was
an inadequate performance measure. V17 admitted 68 MRK samples and one
456-equation low-grade OPE block during discovery; the latter supplied
390 independent restrictions in 489.729 seconds. Its discovery candidate
has 22032 directions instead of 22489, but 400158698 nonzeros instead of
384425474. The reported refinement time rose from 4600.320 to 5121.300
seconds. After the full collinear fit, however, V17 has 2984 directions
and 1082630 nonzeros, compared with V16's 3056 directions and 5907390
nonzeros. Its later MRK fit removes no further directions. Both schedules
reach 2877 directions and 1645065 nonzeros after the full OPE fit.

Measured between successive checkpoint timestamps, the V17 collinear,
MRK, OPE and final-seam stages took 3982.214, 124.299, 934.674 and
1251.481 seconds, respectively; these intervals include checkpoint work.
The corresponding V16 intervals are 4183.009, 297.208, 964.827 and
1309.844 seconds. Thus downstream savings offset the larger discovery
cost in this run. Overall wall time was only 84.241 seconds, or 0.66\%,
shorter, while CPU time increased by about 3.45\% and peak RSS increased.
These are operational measurements at different primes and with different
background workloads, so they establish comparable observed wall times,
not a controlled speedup. V17 remains selected for its validated admission
schedule and acceptable resource use; no global performance optimum is
claimed.

\subsection{Final timings, recovery and resource use}
The eight completed native primes and their conversions have the following
measured wall times. GB denotes $10^9$ bytes, not GiB. The first row is the
V16 pilot; all remaining native rows used V17.
\begin{center}
\begin{tabular}{rrrr}\toprule
Prime index & Native seconds & Native peak RSS (GB) & Shift seconds\\\midrule
1 & 12711.891 & 34.628 & 64.061\\
2 & 12627.650 & 38.581 & 64.129\\
3 & 12859.913 & 38.594 & 67.581\\
4 & 13382.666 & 39.141 & 69.883\\
5 & 13030.380 & 38.602 & 68.881\\
6 & 12911.333 & 38.513 & 69.589\\
7 & 12781.880 & 38.711 & 64.767\\
8 & 13349.812 & 38.582 & 68.285\\\bottomrule
\end{tabular}
\end{center}
All eight have zero surviving directions after their complete original
native joining sweeps and pass their fresh supplied collinear, MRK and OPE
checks. V17 successfully admits the previously postponed MRK/OPE discovery
work, but the completed measurements do not show a large overall ten-loop
speedup over V16. More early physical pivots can increase discovery fill.
The earlier controls, profiling and unsuccessful schedules remain relevant;
no claim of a globally optimal algorithm or engineering configuration is
made. In this final stage, asymmetric reconstruction removes the need to
finish further expensive native primes solely to satisfy a symmetric
coefficient bound.

An abrupt host reboot on 8 October interrupted the transition from native
prime 7 to its shift. System logs do not establish the cause: they contain
no recorded orderly shutdown, nearby OOM, thermal shutdown or kernel panic.
The earlier unrelated OOM is not evidence about this reboot. Native prime 7
was not recomputed. Its bytes matched the durable ledger and every
coefficient of its final sparse checkpoint. Empty checksum and measurement
metadata were restored from that independent evidence with the empty
originals retained. A first resumed shift exited before calculation because
the runtime libraries formerly under \texttt{/tmp} had disappeared.

A permanent, hash-pinned runtime was prepared without changing the frozen
executables. Complete eight-loop native, nine-loop native and nine-loop
shifted controls reproduced every reference byte before continuation.
The resumed seventh shift then completed normally. The replacement user
service guarded runtime hashes at startup and retained the 48 GB pressure
threshold, 50 GB hard limit, zero swap and 16 internal threads. Future
producers should also fsync payloads, sidecars and parent directories before
committing a ledger: atomic rename alone does not guarantee power-loss
durability. This recovery is recorded separately from the mathematics.

After the eighth prime passed the reserved check, the ninth native solve
was temporarily frozen, preserving its state while exact certification ran.
Its cgroup held 39.708 GB at freezing. The certificate service had an 8 GB
hard limit; with a 0.3 GB reserve the combined bound was 48.008 GB. Both a
controller cleanup handler and the service stop handler automatically thawed
the original job. The 44 origin holdouts completed in 44.998 seconds with a
5.624 GB sampled peak, all agreeing exactly. An independent final review
rehashed 161 artifacts, verified the V16 import, recomputed all 44 origin
predictions and confirmed that they cover every ordered-by-size count
partition of 20 into three nonnegative parts. Only then was the now
redundant ninth native solve stopped and its automatic startup disabled.
Its partial work is not counted as a completed prime.

\subsection{Where the result and certificates are recorded}
The primary native result consists of six unique-origin rational matrices in
\begin{center}
\path{output_hexagon_validation_20261009/native-qq-17p3}.
\end{center}
Its 275246 coordinates contain 246941 nonzero entries. The later native
verification section records its reconstruction, exact original-seam and
origin certificates, and exact equality with the shifted export.
The equivalent $16+4$ result is the six unique-origin matrices in
\begin{center}
\path{output_hexagon_17p3_20261007/ten-asymmetric-qq-16p4}.
\end{center}
Each filename is \texttt{irrSEW\_16p4\_}\emph{sector}\texttt{\_origin.wxf},
with the sectors and dimensions of the first table. Entries are exact
rationals; the origin matrix has two basis axes, and there are no remaining
homogeneous directions. The \texttt{columns.tsv} file fixes the order of the
518048 invariant coordinates. These matrices use the unchanged FEC16 and
MHV LEC4 chains under \path{output_hexagon_16p4_20261006}; symmetry expansion
uses the same $E$-component metric $(1,3)$ described above. They are a
factorized ordinary symbol, not an expanded list of all length-20 words.

The final receipt and its independent audit are
\begin{center}
\path{output_hexagon_17p3_20261007/ten-asymmetric-runner/final-verification.json},\\
\path{output_hexagon_17p3_20261007/ten-asymmetric-runner/independent-final-review.json}.
\end{center}
The QQ manifest records the source chains, all eight prime producers, the
frozen candidate, known-result controls and input hashes. The native and
shifted V16/V17 chains remain preserved. The rational candidate contains
491197 nonzero entries, with maximum numerator and denominator sizes
44 and 231 bits. Its independent certificates establish exact joining and
44 exact origin constraints. The other physical validation is the stated
finite supplied COL/MRK/OPE coverage. The empirical ten-loop bonus
restriction remains explicit; neither these certificates nor the observed
finite-field uniqueness proves unrestricted uniqueness or determines
function-level constants invisible to ordinary symbols.

The earlier attempts below explain the engineering changes and remain part
of the reproducibility record.
The FEC15 and FEC16 input generation completed in 422.028 seconds with a
13.702 GB peak. FEC16 has 63882715 compact nonzeros and 232-bit coefficient
height. The V7 eight-loop control and allocator probes have passed.
The first V7 ten-loop trial used a 16 GiB eviction trigger. It was manually
superseded after 1670.57 seconds, with a 26.15 GB peak, twelve completed
evictions and last reported rank 274251. Its counters at 1588 seconds showed
about 56.5 GB read and 238.5 GB written, compared with 29.6 GB of allocated
scratch-file storage. The subsequent V7 retry raised the trigger to 32 GiB to
reduce that page turnover. It used the same frozen executable, mathematical
constraints and 50 GB RSS guard. The trigger changes storage policy only.
The last recorded state of the V7 retry has rank 317799, with 2783740982
retained nonzeros and a sampled peak RSS of 49.6386 GB. The process is no
longer present, and its measurement manager left no final shutdown report.
The interruption cause is not established. It produced no completed native
prime or candidate checkpoint; the inputs and logs are preserved.

V8 passed its storage and reducer controls. It adds a continuous
background pager, since allocator-only checks leave long elimination and
short-relation sweeps unchecked between allocations. A 1100 MiB allocator
probe passes concurrent writes and an automatic eviction with no allocation
callbacks, preserving every tested coefficient byte.

V8 also removes redundant writes from short-relation propagation. The old
pass assigns each retained coefficient back to itself even when its
substitution is unchanged; with billions of nonzeros this dirties a large
fraction of the file-backed basis. The new pass writes only an actual
substitution or compaction. Furthermore, new triangular rows contain no
earlier pivot columns. If every new row has more than two entries, old short
relations cannot reduce them, and no new singleton or doubleton relation is
introduced. In that case the previous substitution state remains exactly
valid and the full sweep is a no-op. When a new short row exists, the
complete reverse propagation still runs. The legacy-reducer comparison passed for singleton, doubleton, cascading
zero and dense no-op cases at both a small prime and the production prime.
The $12+4$ eight-loop control passed all sewing and physical checks and
matched the known coefficient payload exactly; it took 134.344 seconds,
peaked at 2.209 GB and completed 39 evictions.

V9 retained these mathematical and storage kernels and added explicit
selection of the native prime. This allows one prime per process, with
conversion and reconstruction between solves. V10 additionally admits direct physical rows during discovery, as described
above. Its $17+3$ scheduling target is 9000 directions. The managed
user service enforces
a 48 GB memory-pressure threshold and a 50 GB hard cgroup limit, alongside
the process-tree RSS and system-reserve guards. Thus a chat interruption
need not terminate the numerical process. Each successful prime retains
its native cache and sparse checkpoints; only its temporary mapped working
file is eligible for removal after the process exits.

The completed rational candidate and exact certificates now supplement the
one-prime dimension chains. These separate results establish rational
joining and the reserved origin predictions; the finite-field dimension
chains alone would not do so.

The experiment preserves earlier stopped attempts, complete input blocks,
private checkpoints, source and executable hashes, and resource reports.
The audit directory is
\texttt{audits/hexagon-ten-loop-2026-10-06}; the numerical experiment is
\texttt{output\_hexagon\_16p4\_20261006}, with the alternative split
recorded in \texttt{output\_hexagon\_17p3\_20261007}. No stopped attempt is labeled a
completed native prime. The completed result has the conditional certification scope stated in the
abstract and in the final verification receipt.


% BEGIN NATIVE BENCHMARK ADDENDUM
\section{The native rational export and reproducible benchmarks}
The native representation deserves its own durable result because the
linear solve takes place in $17+3$ coordinates. It is now stored as six
rational sector matrices under
\begin{center}
\path{output_hexagon_validation_20261009/native-qq-17p3}.
\end{center}
There are 275246 coordinates and 246941 nonzero entries, with no remaining
homogeneous directions in the selected conditional ansatz. The export uses
the original native recurrence identity and preserves the origin-last
convention. Its manifest explicitly identifies an external rational export;
it does not relabel the first prime as a V17 computation.

Seven saved native primes reconstruct the vector with numerator bound
$N=2^{96}-1$ and denominator bound
$D=\lfloor(M-1)/(2N)\rfloor$. Every coefficient agrees at the eighth
saved prime. This is a retrospective check. It is distinct from the two
primes that were genuinely reserved after the earlier shifted candidate
was frozen. An exact rational conversion of the native vector to $16+4$
then reproduces all 518048 coefficients of that independently certified
candidate. Wolfram's independent WXF reader verifies both the native
export against its rational TSV vector and the converted export against
the accepted shifted vector. The conversion checks its exact factorization
identity before writing the output.

The native export also passes every original, unprereduced joining row.
Clearing the recurrence, route and candidate denominators gives an absolute
integer residual bound of 700 bits. Eight saved congruence primes and four
additional residual-only primes give a 721-bit modulus exceeding that
bound. Each of the 2053259 residuals is zero modulo every prime, so it is
zero over $\mathbb Q$. The four extra primes evaluate an existing candidate;
they do not repeat the native solve. All 44 origin predictions also pass
exactly on the native factorization. This checks the native result directly for coordinate and normalization
errors, in addition to its exact relationship to the shifted result.

The exact native seam check took 116.652 seconds with 14.334 GB sampled
peak RSS; the 44 native origin checks took 145.819 seconds with 18.896 GB.
The shifted factorization is cheaper for these checks: its corresponding
costs were 67.667 seconds/4.398 GB and 44.998 seconds/5.624 GB. The primary
native representation can therefore be retained while a proven exact
change of split supplies a cheaper verification representation.

\subsection{Reconstruction should respect coefficient height asymmetry}
The first native reconstruction trial used six primes with a uniform
48-bit numerator bound. It failed to find 31 coordinates and supplied
incorrect candidates for 31 others, which both saved check primes rejected.
That failure is preserved. It does not signal inconsistent modular solutions:
the verified native vector has denominators as large as 315 bits, whereas
the trial's denominator bound has only 312 bits. The coordinates with these
large denominators have particularly small numerators.

For each reduced coefficient $n_j/d_j$, define the product height
$H_j=2|n_j|d_j$. The maximum individual height in this result has 329 bits,
while the product formed from separate global numerator and denominator
maxima has 361 bits. The six-prime modulus itself has 361 bits. Thus a
single rectangular reconstruction bound can miss useful coefficient-wise
asymmetry. A retrospective experiment enumerates the Euclidean convergents
satisfying $2|n|d<M$ and chooses the unique candidate minimizing $|n|d$,
handling zero separately. This is a candidate-selection heuristic when
coefficient bounds are unknown; it is not a substitute for reserved checks.

The experiment reconstructs every native coefficient from the first six
primes, agrees with both remaining saved caches, and reproduces every
coefficient of the known eight- and nine-loop reference exports (69988 and
195155 coordinates). A measured repetition, including those controls,
takes 9.754 seconds and 0.183 GB peak RSS. The initial uninstrumented
experiment took 11.526 seconds. The accepted seven-prime export's provenance
is retained. For future calculations, freezing this reconstruction rule
before two new validation primes would allow a six-plus-two protocol for
both representations on data of this kind. Retrospective success cannot
turn existing primes into fresh checks.

\subsection{Additional physical coverage and phase measurements}
The broader validation is performed on the primary native vector. Exact
conversion transfers its symbol-level conclusions to the shifted
representation. The new checker first reproduces the saved eight- and
nine-loop references: respectively 8/8 collinear equations, 128/144 MRK
equations, and 2100/2618 extended OPE coefficients. A deliberately doubled
eight-loop symbol fails all 64 equations in a separate MRK negative control.
The projection engines themselves are shared with the solver; the new
points and unused physical coefficients provide additional tests, not an
independent implementation of every formula.

The extended ten-loop OPE table contains 3192 coefficients. Its 2280-entry
overlap with the original table agrees exactly. The additional 912 entries
occur at $T S^{11}$, $T S^{13}$, $T^2 S^{12}$ and $T^2 S^{14}$, with all
supplied log powers in the maximal-helicity sectors. None of those entries
was present in the original fit or original holdouts. The original fit did
include solve-role twist-one data, so twist one as a whole is not an unused
holdout. The input plan fixes the additional data, seed and sample offset
before validation of the ten-loop result.

At each of the new primes 1152921504606947147 and 1152921504606947161,
64 collinear projections, 1280 MRK equations and all 3192 extended OPE
coefficients pass without a mismatch. Thus there are 9072 successful
finite-field checks, including the same 912 distinct new OPE coefficients
at each prime. They are validation equations, not 9072 new independent
fitting constraints. The coverage still excludes helicity-zero twist two
and all higher twists. These finite-field tests are separate from the exact
rational seam, origin and split certificates and do not prove every
physical limit as a rational identity.

\begin{center}
\begin{tabular}{lrrrr}\toprule
Check & Prime 1 seconds & Peak GB & Prime 2 seconds & Peak GB\\\midrule
Collinear & 88.187 & 16.962 & 93.250 & 16.438\\
MRK & 275.950 & 16.912 & 288.675 & 16.905\\
Extended OPE & 1157.125 & 18.908 & 1083.189 & 18.995\\
\bottomrule
\end{tabular}
\end{center}
The six ten-loop physical validation jobs take 2986.376 seconds
in sum, with an 18.995 GB sampled peak. Control calculations and exact
certification are accounted for separately.

\subsection{Reproduction scope and remaining algorithmic comparisons}
The benchmark bundle is \path{benchmarks/hexagon-mhv-ten-loop}. Its primary
reference directory is \path{reference/17p3}; the equivalent shifted result
is in \path{reference/16p4}. Coordinate order, recurrence identities, sources,
commands and validation receipts accompany the six sector matrices. The
complete-prime wrapper starts from the prepared FEC and LEC bases in a new
output directory, and reads the reference vector only after the solve for
comparison. Its command construction and input validation have been checked;
a new complete prime has not been run solely to test this wrapper.

This is a sewing benchmark with prepared bases. It does not yet provide a
validated fresh-clone regeneration of all hexagon bases and external physical
inputs. The eight successful native solves took 103655.526 seconds in total
(28.793 hours), with 39.141 GB maximum sampled process-tree RSS. Their combined
CPU time was 1116756.594 seconds, averaging 10.77 cores. The eight modular
conversions took a further 537.176 seconds. These totals exclude basis
preparation, failed or interrupted trials, reboot downtime, reconstruction
and later certification. Memory figures use decimal GB and 100 ms RSS
sampling, within the 50 GB hard cgroup limit.

The measurements identify discovery and collinear fitting as the main
whole-run targets. Structural row and pivot replay across primes could reuse
an elimination order while recomputing every coefficient and right-hand
side. Exceptional pivots must trigger a checked fallback, and the complete
original seam certificate must remain independent. A separate V18 experiment
has now validated final-seam structural replay at eight and nine loops, as
reported below. V26 subsequently tests discovery replay at eight and nine
loops; its ten-loop cost and fully implicit physical projection remain
unmeasured.

The newly checked higher-$S$ OPE data also permit a concrete scheduling
experiment: admit $S^{11}$/$S^{12}$ data earlier while retaining
$S^{13}$/$S^{14}$ as holdouts. The decision should use total subsequent time
and fill, including projection cost, rather than rank gain alone. A weight-dependent triangular log bound is now validated at nine loops by
V22 below, while its ten-loop performance remains unmeasured:
at head weight 17 there are 171 slots, compared with 231 at total weight 20.
Omitted entries and the known eight- and nine-loop results must be checked
before selecting that change.

Finally, the finishing solver currently restricts its affine space after
64 pivots or at the end of a sweep. If fewer than 64 parameters remain,
restricting as soon as the rank equals the remaining dimension could avoid
some work. A subsequent complete zero-residual sweep would still be required.
Existing logs do not locate the point of rank saturation, so this is a
proposal with no measured saving. These remaining comparisons preclude a
claim of global performance optimality.

\subsection{Measured memory-policy improvement in final sewing}
Both replay runs start from the same prime-two post-OPE checkpoint, with
2877 homogeneous directions and 1645065 stored nonzeros including the affine
origin row. They invoke the identical frozen V17 finishing function and
reproduce all 275246 final native coefficients, including all zero entries.
Each completes the full final zero-residual sweep; the comparison changes
allocation policy without changing equations or mathematical scheduling.

\begin{center}
\begin{tabular}{lrrrr}\toprule
Backend & Seam seconds & Process seconds & CPU seconds & Peak GB\\\midrule
Mapped & 1276.920 & 1349.263 & 2781.930 & 13.959\\
Heap & 222.071 & 289.422 & 2268.137 & 14.410\\
\bottomrule
\end{tabular}
\end{center}
The phase wall-time ratio is 5.750, an 82.61\% reduction. Including input
loading, the process ratio is 4.662. The sampled peak increases by only
0.451 GB, and includes loading as well as the finishing phase. Both jobs
use 16 threads, the same O2 binary and the same field prime. They run serially
in mapped-then-heap order with warm input and OS caches on the shared
workstation. This is one trial per backend and does not establish an
end-to-end prime speedup.

The mapped allocator log records at least 8310784 free-range reclamation
calls. Its implementation punches holes in free storage while holding an
allocator mutex. System CPU time is 877.627 seconds for mapped storage and
82.527 seconds for heap storage; average whole-process CPU usage is 2.06
and 7.84 cores respectively. These observations support allocation and
reclamation overhead as a substantial finishing-stage cost. The free-range
byte counters are not measurements of physical disk traffic.

Live cgroup snapshots verify the 48 GB pressure threshold, the 50 GB hard
limit (rounded down to a page boundary), and zero swap allowance. The
observed high/max/OOM counters and memory pressure are zero. Post-exit
\texttt{systemctl} queries returned default values after the transient units
were collected; the saved live snapshots provide the resource-limit evidence.

Heap storage is therefore the measured choice for a separate replay of
this finishing checkpoint. For future complete solves, a checked process
boundary or an ownership-safe policy switch after the large spaces shrink
is worth benchmarking. The result does not support using heap storage
throughout discovery, where memory retention previously mattered. The
benchmark bundle preserves both commands, input and executable hashes,
full logs, resource records and exhaustive coefficient comparisons. The
frozen production results are unchanged.
\subsection{V18 follow-up on eight- and nine-loop references}
The separate V18 package adds three related changes to final sewing. It
contracts the seam factors directly into candidate coordinates, avoiding
construction, sorting and merging of each full source row. After an affine
restriction it continues at the next block: a previously satisfied row
remains zero under the common left transformation of directions and origin.
Finally, it stores independent source-row addresses for later primes,
regenerates every modular coefficient and right-hand side, solves the small
system and lifts once. Layout, identity and rank checks guard reuse. Every
route retains a complete final operator check. The earlier discovery and
physical scheduling, CG conventions and bonus ansatz are unchanged.

The saved validation has 19 stages and a 208-file build receipt. An additional
review verifies their hashes, reported measurements and every first-prime
native coefficient against the original saved payload. Both new primes pass
the current physical holdouts and their converted vectors agree with the saved reference values. The original
comparator could select stored finite-field caches; those calls therefore
did not establish fresh reduction of the rational reference files. The V22
follow-up below verifies forced rational reduction at nine loops. Phase controls separately replay
all 133214 and 372388 factored sewing rows through the original implementation.
These are finite-field lower-loop controls, not new rational reconstructions
or replacements for the ten-loop exact certificates.

\begin{center}
\begin{tabular}{lrrrr}\toprule
Run & V17 seconds & V18 seconds & V17 peak GB & V18 peak GB\\\midrule
Eight loops, 2 GiB cache & 60.861 & 58.907 & 1.871 & 1.874\\
Nine loops, 2 GiB cache & 511.502 & 495.895 & 5.957 & 5.950\\
Nine loops, V18 3 GiB cache & 511.502 & 467.490 & 5.957 & 5.984\\
\bottomrule
\end{tabular}
\end{center}
With equal cache settings, complete-prime wall time falls by 3.21\% at eight
loops and 3.05\% at nine loops. The 8.60\% reduction in the last row combines
algorithm and cache changes. Each nine-loop terminal OPE head occupies
876573360 bytes; three require 2629720080 bytes, exceeding 2 GiB but fitting
within 3 GiB. The enlarged cache reduces terminal generations from four to
three and removes about 29.9 seconds of logged head-generation work. Peak
memory is essentially unchanged in these trials; this is not a general
memory guarantee.

At the same nine-loop finishing checkpoint, V17 takes 22.075 seconds, V18
learning takes 6.989 seconds, and replay takes 3.391 seconds. The actual later
prime reuses 380 rows and one lift: the small solve takes 0.040 seconds and
complete verification takes 3.337 seconds. At eight loops the checkpoint
already has zero directions, so its improvement concerns verification, not
learning constraints.

These measurements use matched O3 builds, heap allocation, 16 threads and
prepared bases, with one complete-prime trial per configuration. Numerical
timings begin after the ten-loop allocator comparison. The jobs have a 20 GB
hard cgroup cap, no swap and zero recorded limit/OOM events. The ten-loop
allocator comparison used O2, and these lower-loop factors cannot be
multiplied into a ten-loop speedup estimate. No ten-loop V18 production
selection is made. A bounded ten-loop pilot, repeated timings and structural
ten-loop discovery replay remain useful next comparisons. V26 below
subsequently tests discovery replay at eight and nine loops.

The source package and detailed report are in
\path{output_hexagon_v18_20261009}; the benchmark archives the reviewed
receipts under \path{evidence/v18-lower-loop}. The accepted ten-loop
coefficients, checks and producer identities remain unchanged.

\subsection{V22: compressing physical operators and intermediate OPE storage}
The historical nine-loop study measures complete process user-plus-system
CPU time. Its then-applied gate requires more than 15\% improvement over original V17,
with no more than 5\% higher sampled peak RSS, using the smaller V17 peak
as the conservative reference. Independent review recomputes the means,
checks 14 run receipts and 834 source/build file bindings, and repeats all
native coefficient-payload comparisons.
This historical V17 comparison does not establish the current V16-baseline
qualification clarified below.

\begin{center}
\begin{tabular}{lrrrr}\toprule
Configuration & Trials & CPU-hours & Wall seconds & Peak GB range\\\midrule
V17 / 16 threads & 2 & 1.481645 & 508.864 & 5.903--6.011\\
V20 / 16 threads & 2 & 1.049393 & 384.887 & 5.877--5.935\\
V21 / 16 threads & 3 & 1.039490 & 380.377 & 5.849--5.939\\
V21 / 8 threads & 1 & 1.178127 & 684.994 & 5.323\\
V22 / 16 threads & 2 & 0.918766 & 347.391 & 5.840--5.890\\
\bottomrule
\end{tabular}
\end{center}
CPU and wall means use first-prime repetitions; RSS ranges include every
field and repetition. V22 reduces CPU time by 37.99\% relative to V17 and
11.61\% relative to V21 with 16 threads. The eight-thread control loses on
CPU and wall time. Its wall-time-driven physical schedule changes the
discovered space, so it is a configuration comparison rather than pure
hardware scaling. The initial V19 candidate saved only 11.67\% against the
final two-trial baseline and failed the gate.

The main change compresses the affine collinear suffix operator together
with every sampled right-hand-side column. Write the augmented operator as
$Q=[T\mid b_1\ \cdots\ b_m]$. Selecting original rows $I$ and explicitly
verifying a certificate $Q=CQ_I$ preserves the full row span, including
RHS-only equations. The observed suffix count falls from 36 to 14. Rechecking
exact zero-head support for each field and sample batch reduces it to nine.
The smaller operator avoids repeated large projected sparse products while
retaining the scalar physical holdouts and complete seam verification.

V22 then stores only logarithm powers $a+b\leq w$ at intermediate head
weight $w$, retaining all $T/S$ jet coefficients and restoring the existing
external layout before projection and cache use. At weight 15, 136 slots
replace 190. The implementation passes 2160 square-reference operator cases
and all 1178100 coefficients of complete FEC8 heads at each of two primes.
Mean OPE CPU time decreases from 1313.38 seconds in V21 to 1000.55 seconds
in V22. Collinear timing also improves after the shared recurrence gains
separate input/output strides. No isolated attribution experiment was made,
so the full improvement cannot be assigned solely to OPE storage.

Every one of the 105898 native coordinates agrees at both primes, and the
current collinear, MRK, OPE and complete 372388-row factored seam checks pass.
An additional reference audit corrects the earlier comparator ambiguity:
cache-free snapshots match all original rational reference files by hash,
forcing genuine rational reduction. Both V20 shifted outputs agree in all
195155 coordinates; V22 native payload equality with V20 transfers that
check. The older V18 comparison evidence remains preserved with its narrower
scope. This correction does not affect the ten-loop independent Wolfram
readback, exact split equality or original-seam certificates.

These are serialized sewing trials on prepared bases, with a 20 GB hard
cgroup limit, no swap and zero recorded limit/OOM events. Physical cache is
4 GiB and OPE cache is 3 GiB, except original V17 uses 2 GiB. There is no
confidence interval, cold-basis timing or V22 ten-loop measurement. The
50 GB production limit is unchanged. Discovery still costs about 1033 CPU
seconds per nine-loop run, motivating guarded row/pivot replay and
precomposition of the double-bonus tail basis as further comparisons.

\subsection{Additional triple-entry checks and their limits}
The separate mathematical audit checks the complete 21-term Eq.~(92) of
ref.~\cite{dixonliu} and both cyclic images. The completed native ten-loop
symbol passes at two primes after fitting. These are additional modular
holdouts, not a new rational triple certificate or an all-loop proof.
The constraints were not used in V22 or in the accepted ten-loop solve.

On the audited local source they add rank three beyond the double bonuses,
reducing 65 local triple directions to 62. Rational local witnesses prove
that the certified local source alone does not imply these relations;
without a global physical-prefix extension they are not counterexamples to
the amplitude relations. The initial truncated transcription omitted
continuation lines and was invalid. Its apparent failures are superseded
by the complete v3 checks, including exact known eight- and nine-loop tests.
A separate quadruple calibration returned 169 rather than 166 and remains
unresolved; it supplies no amplitude-error claim or solver constraint.

The numerical evidence and pinned nine-loop profile are archived under
\path{evidence/v22-nine-loop} in the benchmark. The mathematical source and
receipts are in \path{bonusrelation/solver_audit_20261009}. All claims of
uniqueness retain the stated conditional ordinary-symbol scope.

\subsection{V26: generic discovery replay and triangular affine factors}
The next implementation moves structural reuse into discovery and delays
expansion of triangular affine equations. The core accepts general sparse
finite-field systems: source rows determine ranks, pivots and free variables.
The physical adapter supplies the existing equations and normalization;
no additional bonus assumption or known-answer fitting equation is introduced.

A discovery plan records source positions, pivot groups and bounded
dependency routes, never coefficients from an earlier field. Every matrix
entry and right-hand side is regenerated at the new prime. Discarded rows
are checked against the new basis. An absent dependency route uses ordinary
exact elimination locally, while changed source identity, row streams,
pivots or rank trigger fresh discovery. This preserves the complete system
even on an exceptional prime, though fallback can remove the speed advantage.

The affine factor retains triangular equations
\[
 x_j=b_j-\sum_{k>j}U_{jk}x_k.
\]
Free-variable maps and a shared constant coordinate are composed directly,
avoiding an intermediate fully expanded reduced matrix. The generic API
supports lifting, operator projection and materialization. The current
driver still materializes complete maps at its existing affine checkpoint
boundaries; fully implicit propagation through all physical stages is not
yet a measured implementation.

\begin{center}
\begin{tabular}{lrrrr}\toprule
Configuration & Learn CPU-h & Replay CPU-h & Mean wall s & Peak GB\\\midrule
Eight-loop V26 & 0.097601 & 0.098072 & 47.654 & 1.883\\
Nine-loop fresh V22 & 0.915311 & 0.907920 & 342.590 & 5.949\\
Nine-loop V26 & 0.880155 & 0.864388 & 324.935 & 5.943\\
\bottomrule
\end{tabular}
\end{center}
All three rows use 16 workers. Each mean includes one learning prime and
one replay prime, not statistical repetitions of one prime. The V26
nine-loop mean is 0.872271 CPU-hours: 41.13\% below the historical V17
mean of 1.481645 CPU-hours and 4.32\% below the fresh matched V22 pair
(0.911615 CPU-hours). The eight-loop mean is 0.097837 CPU-hours, a 23.34\%
reduction against its original V17 control. These include learning costs.
The historical V17 comparison changes OPE cache from 2 to 3 GiB as well as
algorithm, and therefore measures a whole configuration.

Here peak RSS is the larger of the 100 ms sample maximum and kernel process
high-water marks recorded by phase instrumentation. The nine-loop value
is 5942951936 bytes, only 0.672\% above the conservative historical sampled
V17 reference of 5903253504 bytes. The eight-loop value is 1883369472 bytes,
0.648\% above its reference. Earlier tables retain their sampled metric.
Thus both tested workloads pass the greater-than-15\% CPU saving and
at-most-5\% RSS growth gate then applied against original V17. The V17
measurements retain that identity and do not establish V16 qualification.

The distinction between sampled and kernel peaks changes the engineering
decision. Profiling measured 825 CPU seconds reducing incoming rows and
27 seconds searching for new pivots, motivating dependency-route replay.
A 256 MiB route cache passed correctness but reached 2.186 GB at eight
loops. Reclaiming unused allocator pages reduced its kernel peak to
2.060 GB, still too high; the sampler alone reported 1.884 GB and missed
this brief allocation spike. V26 uses a fixed 64 MiB route budget for every
input, with reclamation only after temporary owners are destroyed at joined
phase boundaries. V24/V25 have slightly smaller nine-loop CPU means but
fail this cross-instance memory criterion. Their small time differences do
not establish a general performance ordering.

Independent review verifies 21 native run receipts, all 221 selected-build
file hashes, every native payload comparison, and the recomputed CPU/RSS
gate. Generic controls cover 96 field cases, 48 structural replays and
72 affine-map comparisons against independent dense elimination, including
rank-deficient and inconsistent systems and two exceptional-prime rank
changes. Deliberately absent routes pass 48 fallback controls. Five other
matrix families agree with both the previous implementation and all original
equations. The selected runs reproduce all 38989 eight-loop and 105898
nine-loop native coordinates at both fields; complete sewing sweeps and
fresh collinear, MRK and OPE checks pass. The nine-loop rational reference
chain retains the earlier forced-QQ-only reduction audit.

Runs were serialized with native heap storage, 4 GiB physical cache,
3 GiB OPE cache, an 18 GB pressure threshold, a 20 GB hard cap and no swap.
Time includes input loading, discovery, fitting, checkpoint and plan I/O,
and final checks. Compilation, cold basis/RHS generation and a new rational
reconstruction are excluded. All 21 fresh native runs together consume
10.688 CPU-hours. These controls establish correctness and measured costs
on the tested workloads, not universal speedups or global optimality.

This study motivates compressed dependency routes, selection by work saved
per byte, a deterministic work-based physical schedule, and implicit affine
factors across physical projection stages. V27 below tests the first three;
complete implicit projection remains unmeasured. A time-based
schedule can change the source-row stream and safely force rediscovery.
V26 ten-loop performance remains unmeasured; lower-loop savings cannot be
multiplied into the accepted ten-loop timings. The 50 GB production limit,
native and shifted rational references, and original producers are preserved.
The independently reviewed receipts are in \path{evidence/v26-generic};
full packages and plans remain in \path{output_hexagon_v23_20261009}.

\subsection{V27: further experiments and the corrected working selection}
V27D with 16 workers is the current working nine-loop benchmark. The
qualification target remains strictly more than 15\% CPU saving against
V16, with at most 5\% higher RSS; V26 supplies an incremental comparison.
The earlier recommendation to retain V26 incorrectly reset the qualification
threshold against V26. There was no separate correctness failure motivating
that recommendation. Direct nine-loop qualification against V16 remains
unverified: the historical comparison measured V17 and must retain that
identity. Working selection does not assert that the missing V16 comparison
has passed or that ten-loop V27 performance has been measured.

\begin{center}
\begin{tabular}{lrrr}\toprule
Complete native workload & Mean CPU-h & Mean wall s & Peak GB\\\midrule
Nine-loop V26 / 16 workers & 0.874501 & 325.347 & 5.912\\
Nine-loop V27D / 16 workers & 0.867992 & 297.189 & 6.010\\
Nine-loop V27D / 8 workers & 0.838348 & 468.654 & 5.374\\
Eight-loop V26 / 16 workers & 0.099578 & 48.714 & 1.902\\
Eight-loop V27D / 16 workers & 0.063614 & 26.606 & 1.784\\
\bottomrule
\end{tabular}
\end{center}
Each row averages a learning prime and a replay prime, not timing
repetitions. Peak memory is the maximum of the 100 ms RSS sample and kernel
process high-water marks. V27D with 16 workers saves 0.74\% CPU time and
8.65\% wall time, and raises peak RSS by 1.65\% against fresh V26. The small
CPU difference is not an established statistical speedup. Eight workers save 4.13\% CPU
against that 16-worker baseline but take longer. A separate matched
eight-worker first-prime comparison gives 0.869639 versus 0.845487 CPU-hours,
a 2.78\% saving. The eight-loop 36.12\% saving against V26 does not establish
nine-loop qualification against V16. A historical eight-loop V16 control
gives 49.32\% CPU saving and 1.44\% RSS growth for V27D, but compares one
original mapped/cache-configured V16 prime with a new heap learning/replay
pair. It is not a matched nine-loop qualification experiment.

The implementation tests four reusable directions. Ordered bit queues,
balanced row ranges, reusable worker scratch and exact density-based
sparse/dense dispatch improve the numerical kernels. Delayed 128-bit
accumulation corrects arithmetic overflow exactly; fields larger than
61 bits retain canonical reduction. Scheduling compares physical rank
gain per arithmetic work against seam work, including retained fill, and
records task decisions for replay. An augmented feature-space compressor
certifies discarded equations using all RHS and origin columns. Finally,
delta-varint dependency routes are selected by estimated updates saved per
byte. Their accounting budget is 64 MiB, with an 8 MiB logical learning
budget; incomplete routes finish with exact residual elimination. These
are transformations of the supplied system, not additional physical or
bonus constraints.

First-prime ablations at nine loops save 3.25\% CPU for the kernel bundle,
1.43\% for the compression extension and 0.64\% for packed routes. The new
schedule costs 0.13\% more CPU in its local comparison. Packed routes save
1.65\% over learning plus replay relative to the legacy route cache.
These are interacting measurements with other features enabled, not additive
gains or confidence intervals. Disabling packed routes restores legacy
structural replay rather than removing replay altogether.

Failures supply useful engineering constraints. Uniform sampled density
aliased polynomial layouts, requiring exact density inspection. The initial
work-budget-only physical schedule admitted expensive MRK/OPE work too
early: V27B averaged 1.103377 CPU-hours per nine-loop prime. Comparing
physical gain/work against seam work recovered most of that regression.
The MRK feature-space operator remained full rank, so the new compression
adapter removed no rows there. All completed variants stayed numerically
correct; their separate source packages, executables and logs are preserved.
The scheduling regression concerns early execution, not availability of
physical constraints. All required equations or lazy generators should be
registered from the start; expensive blocks may remain pending according
to expected rank gain, total arithmetic cost, fill and memory. Deferred
blocks must preserve the affine RHS and origin under coordinate changes
and remain subject to completion accounting and final checks.

The phase data distinguish elapsed time from CPU cost. First-prime discovery
uses 930.056 CPU seconds in 126.644 wall seconds for V26 with 16 workers,
1112.650 CPU seconds in 104.756 wall seconds for V27D with 16 workers, and
868.333 CPU seconds in 140.247 wall seconds for V27D with eight workers.
Lower worker count also improves MRK CPU time while worsening collinear
and OPE CPU time. Phase-dependent worker selection is an untested next
comparison. Hardware counters were unavailable, so these timings do not
establish cache-miss or memory-bandwidth rates. Nested phase records must
not be summed as independent costs.

Independent review checks all 27 native runs, 909 source/build file bindings
across four packages, serial timing intervals, all native coefficients and
the recalculated measurements and ablations. Every run reproduces the
38989 eight-loop or 105898 nine-loop reference coordinates and passes
complete sewing and fresh collinear, MRK and OPE checks. The prior
forced-QQ-only reference chain is retained. New and legacy switches pass
96 field cases, 48 replays, two exceptional-prime fallbacks and 72 affine-map
comparisons. Additional controls cover 144 augmented compression cases,
64 deterministic schedules, 756718 queue operations, route damage/budget
guards, exact density and persistent cycles, and 5077 deliberately exercised
overflow corrections against canonical arithmetic.

All 27 native trials together consume 16.098 CPU-hours. They run serially
with prepared inputs, heap allocation, 4 GiB physical cache and 3 GiB OPE
cache, under a 20 GB hard limit, 18 GB pressure threshold and no swap.
Recorded cgroup events are zero. Timings include loading, checkpoint and
plan I/O and final checks, but exclude compilation, cold inputs and a new
rational reconstruction. The review is archived in
\path{evidence/v27-experimental}; complete sources and plans remain in
\path{output_hexagon_v27_20261009}. The current working selection and its
independent amendment review are in \path{evidence/v27-selection-amendment}.
Original reports and receipt reviews retain their old selection decision as
historical evidence; the amendment supersedes that recommendation. The old
execution gate is preserved byte for byte with an explicit verified
relocation. Legacy comparison flags refer to the superseded V26 threshold,
not the unverified V16 qualification. At this selection stage ten-loop V27
performance was unmeasured; the later first-prime result follows below.
The accepted ten-loop producers, 50 GB production cap and
conditional uniqueness scope are preserved.

\subsection{First fresh V27D ten-loop prime and the memory regression}
The fresh V27D solve at $p=1152921504606847009$ now reproduces every one
of the 275246 native and 518048 shifted coefficients of the preserved
modular result. Independent receipt review checks the complete 989800-row
factored seam and the fresh supplied collinear, MRK and OPE checks over
this field. These are modular checks, distinct from the original unreduced
QQ sewing certificates of the already certified result. The new run's
rational reconstruction and exact certificates remain pending; its second
prime was running at the first-prime snapshot. Its completed replay
measurement is reported in the following subsection.

\begin{center}
\begin{tabular}{lrrr}\toprule
Native learning prime & CPU-hours & Wall seconds & Sampled peak GB\\\midrule
Historical V16 / 16 workers & 37.285539 & 12711.891 & 34.628\\
Fresh V27D / 16 workers & 23.379927 & 6599.910 & 46.188\\
\bottomrule
\end{tabular}
\end{center}
The observed savings are 37.295\% in CPU time and 48.081\% in wall time,
but sampled peak RSS grows by 33.383\%. Thus the CPU threshold passes
while the 5\% memory tolerance fails: this ten-loop configuration does not
pass the combined criterion. Staying within the absolute 50 GB ceiling
is a separate property. The V27D kernel process high-water mark is
46.208 GB; the table compares sampled RSS on both sides. The recorded
service-wide peak was 47.774 GB, with zero high/max/OOM events at the
comparison snapshot.

V16 used mapped storage with an 8 GiB OPE-head cache, while V27D uses
heap storage with a 3 GiB OPE-head cache. Both use 16 workers and prepared
basis inputs. This is one historical configuration comparison, not isolated
algorithm attribution, statistical repetition or nine-loop V16 qualification.
At that stage replay and total reconstruction costs remained unmeasured. Conversion
adds 65.740 wall seconds and 0.031945 CPU-hours, at 10.785 GB sampled peak
RSS, and is excluded from the native comparison.

\begin{center}
\begin{tabular}{lrr}\toprule
Native phase & CPU-hours & Wall seconds\\\midrule
Discovery & 13.395889 & 3691.310\\
Collinear restriction & 4.871917 & 1266.440\\
MRK restriction & 0.479636 & 125.562\\
Low-$S$ OPE restriction & 2.555983 & 809.382\\
Complete final sewing & 0.092571 & 42.516\\
\bottomrule
\end{tabular}
\end{center}
Discovery and collinear restriction account for approximately 57.3\% and
20.8\% of the whole-native CPU cost. Loading, checks and other overheads
are not all represented in this table, and nested profile records must not
be added again. The final MRK restriction gains no rank in this run.
Skipping redundant elimination would require a reusable exact span or
equation certificate; a zero-rank observation alone does not authorize
dropping its required equations or independent final checks.

The peak identifies a concrete generic memory hypothesis. Discovery retains
1649199276 entries, or 19.790 GB of row payload, before
\path{streaming_echelon.hpp} calls \path{triangular_solver_2}.
In the frozen \path{SparseRREF/sparse_mat.h}, the transpose builder appends
row incidence for every stored column with mutex-protected
\texttt{push\_back}. Four-byte indices alone require at least 6.597 GB for
this full transpose, before capacity slack and metadata. This lower bound
does not estimate the savings from removing any particular incidence.
The inspected recursive and scalar backsolve consumers access pivot columns.

A candidate replacement would count and allocate exact-capacity incidence
only for required pivot columns. Before selection, all consumer accesses
must be proved covered and the new fill/indexing path tested on generic
finite-field, rational, affine, rank-deficient and exceptional-prime systems.
Actual pivot incidence, live transpose storage, affine materialization and
allocator retention must be measured separately. A 20-second backsolve
sample used approximately 15.06 active-core equivalents, so the observation
does not diagnose insufficient task parallelism.

The proposal remains unimplemented and unmeasured. The review changed no
active solver, controller, inputs, checkpoints or completed coefficients and
launched no competing heavy benchmark. The replay run retains 48 GB
MemoryHigh, 50 GB MemoryMax and no swap. Evidence is archived in
\path{evidence/v27-ten-prime1}, with active outputs under
\path{output_hexagon_ten_v27d_20261009}. The original certified references
and conditional bonus assumptions are preserved.

\subsection{Second fresh V27D prime: the scope of replay savings}
The solve at $p=1152921504606847067$ passes the recorded complete
989800-row factored seam and fresh supplied collinear, MRK and OPE checks.
Independent review again compares all 275246 native and 518048 shifted
coefficients with their preserved same-prime references. Prime three is
active; this fresh chain's rational reconstruction, new-prime validation
and final exact certificates remain pending.

\begin{center}
\begin{tabular}{lrrr}\toprule
Native V27D / 16 workers & CPU-hours & Wall seconds & Sampled peak GB\\\midrule
Prime one, learning & 23.379927 & 6599.910 & 46.188\\
Prime two, replay & 23.329189 & 6540.289 & 47.601\\
\bottomrule
\end{tabular}
\end{center}
Observed whole-native CPU decreases by 0.2170\%, wall time by 0.9034\%,
and sampled peak RSS increases by 3.059\%. These are different primes,
without timing repetitions or a same-prime replay-off control, so they do
not isolate causal replay gains. Native solves plus both conversions total
46.770291 CPU-hours. The second conversion adds 64.429 wall seconds and
0.029231 CPU-hours; a complete reconstruction cost is not yet available.

\begin{center}
\begin{tabular}{lrrrr}\toprule
Phase & Learn CPU-h & Replay CPU-h & Learn wall s & Replay wall s\\\midrule
Discovery & 13.395889 & 13.385500 & 3691.310 & 3655.980\\
Final sewing & 0.092571 & 0.046964 & 42.516 & 21.163\\
\bottomrule
\end{tabular}
\end{center}
Final-seam replay reduces contractions from 989800 to 2877 and affine
lifts from 34 to one, while retaining the full 989800-row verification.
The phase CPU reduction is 49.27\%, but its difference is only 0.1951\%
of the first prime's whole-native CPU cost. The replay solve takes
0.558 wall seconds; full verification still takes about 20.6 seconds.
Thus an effective optimization of this phase has little effect on the
dominant whole-solve cost.

Discovery CPU falls by only 0.0776\%. Learning performs
12329356644947 coefficient updates and replay 12332936520495, both about
12.33 trillion. Recorded route steps are 64870291 alongside 3231942664
ordinary queue steps, or 1.968\% of their combined count. This is step
coverage, not avoided arithmetic or saved CPU. There are 3739 local residual
completions among 4402 route uses. The inspected implementation follows
each routed prefix by exact elimination of all remaining pivots; these
completions are not global replay failure. The profile records structural
replay for this completed discovery.

At the comparison snapshot the service records 4481 MemoryHigh events,
48.019 GB cgroup peak, and zero MemoryMax or OOM events. Event counts are
cumulative service observations, not stall durations or isolated phase
costs. Process sampled RSS peaks at 47.601 GB, also the rounded kernel
high-water mark. Limits remain 48 GB soft, 50 GB hard and no swap. This
pressure evidence does not repair the first-prime relative RSS gate failure.

The same-prime historical producer was V17, with 38.571675 CPU-hours and
38.581 GB sampled peak RSS. V27D uses 39.52\% less CPU and 23.38\% more
RSS in that historical configuration comparison. The V17 identity must
not be changed to V16; fixed V16 qualification and the unverified direct
nine-loop comparison retain their stated scope.

The profiling priorities are discovery update arithmetic and locality,
backsolve transpose storage and allocation lifetimes. Counts alone do not
separate arithmetic from bandwidth limits. Route admission should use net
queue, decoding and residual work avoided, rather than coverage alone.
The pivot-column-only exact-capacity transpose still requires independent
generic validation. No active producer change, enlarged route cache or
competing heavy benchmark was introduced. Required physical equations,
affine RHS/origin and final checks are preserved. The independent review
and receipts are in \path{evidence/v27-ten-prime2}.

% END NATIVE BENCHMARK ADDENDUM

\begin{thebibliography}{9}
\bibitem{dixonliu}
L. J. Dixon and Y.-T. Liu,
\emph{An Eight Loop Amplitude via Antipodal Duality},
\href{https://arxiv.org/abs/2308.08199}{arXiv:2308.08199}.
\bibitem{flintmemory}
FLINT developers, \emph{Memory management},
\href{https://flintlib.org/doc/memory.html}{FLINT documentation};
installed-runtime ownership implementation checked against
\href{https://github.com/flintlib/flint/blob/v3.2.2/src/fmpz/link/fmpz_single.c}{FLINT 3.2.2 integer cache source}.
\end{thebibliography}
\end{document}
